\section{Ocena polityki}

\begin{frame}{Wyprowadzenie wzoru : Krok 1 -- Definicja z twierdzenia}
\frametitle{Wyprowadzenie wzoru : Punkt wyjścia}
Z twierdzenia o redukcji polityk:
\begin{equation}
J^{\alpha, \beta}_{\mu, \phi_s}(s) = \sum_a \mu(a\mid s) \left[ r(s, a) + \alpha\beta \sum_{s'} q(s'\mid s, a) \, J^{\beta}_{\phi_s}(s') \right].
\end{equation}
Cel: wyrazić zależność wyłącznie przez $J^{\alpha, \beta}_{\phi_s}(s')$, bez jawnego $J^{\beta}$.
\end{frame}

\begin{frame}{Wyprowadzenie wzoru : Krok 2 -- Definicje wartości stacjonarnej}
\frametitle{Wyprowadzenie wzoru : Definicje dla $\phi_s$}
Dyskontowanie wykładnicze:
\begin{equation}
J^{\beta}_{\phi_s}(s') = \mathbb{E}_{\phi_s} [r(s', \cdot)] + \beta \mathbb{E}_{\phi_s, q} [J^{\beta}_{\phi_s}(s'') \mid s'],
\end{equation}
gdzie $\mathbb{E}_{\phi_s} [r(s', \cdot)] = \sum_{a'} \phi_s(a'\mid s') \, r(s', a')$.

Quasi-hiperboliczne:
\begin{equation}
J_{\phi_s}(s') = \mathbb{E}_{\phi_s} [r(s', \cdot)] + \alpha \beta \mathbb{E}_{\phi_s, q} [J^{\beta}_{\phi_s}(s'') \mid s'].
\end{equation}
\end{frame}

\begin{frame}{Wyprowadzenie wzoru : Krok 3 -- Izolacja $\alpha \beta J^{\beta}_{\phi_s}(s')$}
\frametitle{Wyprowadzenie wzoru : Izolacja $\alpha \beta J^{\beta}_{\phi_s}(s')$}
Pomnóż (2) przez $\alpha \beta$:
\begin{equation}
\alpha \beta J^{\beta}_{\phi_s}(s') = \alpha \beta \mathbb{E}_{\phi_s} [r(s', \cdot)] + \alpha \beta^2 \mathbb{E}_{\phi_s, q} [J^{\beta}_{\phi_s}(s'') \mid s'].
\end{equation}

Pomnóż (3) przez $\beta$:
\begin{equation}
\beta J_{\phi_s}(s') = \beta \mathbb{E}_{\phi_s} [r(s', \cdot)] + \alpha \beta^2 \mathbb{E}_{\phi_s, q} [J^{\beta}_{\phi_s}(s'') \mid s'].
\end{equation}
\end{frame}

\begin{frame}{Wyprowadzenie wzoru : Krok 4 -- Odejmowanie}
\frametitle{Wyprowadzenie wzoru : Kluczowe odejmowanie}
Odejmij (7) od (6):
\begin{equation}
\alpha \beta J^{\beta}_{\phi_s}(s') - \beta J_{\phi_s}(s') = \alpha \beta \mathbb{E}_{\phi_s} [r(s', \cdot)] - \beta \mathbb{E}_{\phi_s} [r(s', \cdot)].
\end{equation}
Skraca się $\alpha \beta^2 \mathbb{E}[J^{\beta}_{\phi_s}(s'')]$!

Uproszczenie:
\begin{equation}
\alpha \beta J^{\beta}_{\phi_s}(s') - \beta J_{\phi_s}(s') = - (1 - \alpha) \beta \sum_{a'} \phi_s(a'\mid s') r(s', a').
\end{equation}
Przekształcenie:
\begin{equation}
\alpha \beta J^{\beta}_{\phi_s}(s') = \beta J_{\phi_s}(s') - (1 - \alpha) \beta \mathbb{E}_{\phi_s} [r(s', \cdot)].
\end{equation}
\end{frame}

\begin{frame}{Wyprowadzenie wzoru : Krok 5 -- Podstawienie}
\frametitle{Wyprowadzenie wzoru : Podstawienie do (3)}
Podstaw (10) do wzoru początkowego:
\begin{align}
J^{\alpha, \beta}_{\mu, \phi_s}(s) &= \sum_a \mu(a\mid s) \left[ r(s, a) + \sum_{s'} q(s'\mid s, a) \left( \beta J_{\phi_s}(s') - (1 - \alpha) \beta \mathbb{E}_{\phi_s} [r(s', \cdot)] \right) \right] \\
&= \sum_a \mu(a\mid s) \left[ r(s, a) + \sum_{s'} q(s'\mid s, a) \left[ -(1 - \alpha)\beta \sum_{a'} \phi_s(a'\mid s') r(s', a') + \beta J_{\phi_s}(s') \right] \right].
\end{align}
\end{frame}

\begin{frame}[fragile]{Algorithm 1: Model-Free Policy Evaluation}
\frametitle{Algorithm 1: Model-Free Policy Evaluation for QH}
\begin{algorithm}[H]
\caption{Model-Free Policy Evaluation for QH Discounting}
\begin{algorithmic}[1]
\REQUIRE Stepsizes $(\eta_n)$, $(\theta_n)$; discount factors $\alpha, \beta$; sampling policy $\nu$; evaluation policies $\mu, \phi_s$
\STATE \textbf{Initialize:} $J_0, W_0 \in \mathbb{R}^{|S|}$
\FOR{$n = 0, 1, 2, \dots$}
\FOR{$s \in S$}
\STATE Sample $a \sim \nu(\cdot|s)$
\STATE Observe reward $r(s, a)$ and next state $s' \sim q(\cdot|s, a)$
\STATE Sample $a' \sim \phi_s(\cdot|s')$
\STATE Observe reward $r(s', a')$
\STATE Compute TD target: $r^{\text{target}}_n(s) \leftarrow r(s, a) - (1 - \alpha)\beta \, r(s', a') + \beta W_n(s')$
\STATE Update $W$: $W_{n+1}(s) \leftarrow W_n(s) + \eta_n \left[ \frac{\phi_s(a|s)}{\nu(a|s)} r^{\text{target}}_n(s) - W_n(s) \right]$
\STATE Update $J$: $J_{n+1}(s) \leftarrow J_n(s) + \theta_n \left[ \frac{\mu(a|s)}{\nu(a|s)} r^{\text{target}}_n(s) - J_n(s) \right]$
\ENDFOR
\ENDFOR
\end{algorithmic}
\end{algorithm}
\end{frame}

\begin{frame}{Lemma 1: Operator $T^{\phi_s}$ jako kontrakcja}
\frametitle{Lemma 1: Operator typu Bellmana dla QH}
Definicja operatora $T^{\phi_s}: \mathbb{R}^{|S|} \to \mathbb{R}^{|S|}$:
\begin{equation}
(T^{\phi_s} J)(s) = \sum_a \phi_s(a\mid s) \left[ r(s, a) + \sum_{s'} q(s'\mid s, a) \left[ -(1 - \alpha)\beta \sum_{a'} \phi_s(a'\mid s') r(s', a') + \beta J(s') \right] \right].
\end{equation}
\begin{itemize}
\item $T^{\phi_s}$ jest kontrakcją w normie $\max$ z faktorem $\beta < 1$.
\item Jedyny punkt stały $J = T^{\phi_s} J$ to $J^{\alpha, \beta}_{\phi_s}$.
\end{itemize}
Dowód: $\|T^{\phi_s} J_1 - T^{\phi_s} J_2\|_\infty \leq \beta \|J_1 - J_2\|_\infty$ (twierdzenie Banacha).
\end{frame}

\begin{frame}{Twierdzenie 2: Konwergencja Algorithm 1}
\frametitle{Twierdzenie 2: Konwergencja ocny polityki}
Założenia: kroki uczenia $\eta_n, \theta_n$ spełniają warunki Robbins-Monro: \\ $\sum_n \eta_n = \infty$, $\sum_n \eta_n^2 < \infty$ oraz $\sum_n \theta_n = \infty$, $\sum_n \theta_n^2 < \infty$; \\ dodatkowo $\theta_n / \eta_n \to 0$ (dwie skale czasowe).

Wynik: ciągi $(W_n, J_n)$ zbieżne prawie na pewno do $(J^{\alpha, \beta}_{\phi_s}, J^{\alpha, \beta}_{\mu, \phi_s})$.

\end{frame}